\documentclass[10pt]{article}
\providecommand{\FIGDIR}{fig}
\newcommand{\DIGESTDATE}{15 September 2026}
% ===================== PM Research Watch - shared preamble =====================
\usepackage[T1]{fontenc}
\usepackage[utf8]{inputenc}
\usepackage{amsmath,amssymb}
\usepackage{tgpagella}
\usepackage{tgheros}
\renewcommand{\sfdefault}{qhv}
\renewcommand{\ttdefault}{lmtt}
\usepackage{microtype}
\usepackage[english]{babel}

\usepackage[a4paper,top=17mm,bottom=17mm,left=17mm,right=17mm,
            headsep=6mm,footskip=10mm]{geometry}
\usepackage{graphicx}
\usepackage{xcolor}
\usepackage{colortbl}
\usepackage{tikz}
\usetikzlibrary{calc}
\usepackage[most]{tcolorbox}
\usepackage{titlesec}
\usepackage{fancyhdr}
\usepackage{booktabs}
\usepackage{array}
\usepackage{longtable}
\usepackage{enumitem}
\usepackage{multicol}
\usepackage{ragged2e}
\usepackage[hidelinks]{hyperref}

% ------------------------------------------------------------------ palette
\definecolor{Ink}    {HTML}{14262E}
\definecolor{Deep}   {HTML}{0F4C5C}
\definecolor{Teal}   {HTML}{1F7A8C}
\definecolor{Sky}    {HTML}{4FA3C4}
\definecolor{Sage}   {HTML}{6FA287}
\definecolor{Amber}  {HTML}{E3A008}
\definecolor{Clay}   {HTML}{B5643C}
\definecolor{Coral}  {HTML}{D1495B}
\definecolor{Violet} {HTML}{6A5B9E}
\definecolor{Slate}  {HTML}{4A6572}
\definecolor{Paper}  {HTML}{FBF9F5}
\definecolor{Mist}   {HTML}{EEEAE2}
\definecolor{Rule}   {HTML}{D9D3C9}

\pagecolor{Paper}
\color{Ink}

\hypersetup{colorlinks=true, linkcolor=Deep, urlcolor=Teal, citecolor=Teal,
            pdftitle={PM Research Watch}, pdfauthor={Automated daily digest}}

% ------------------------------------------------------------------ headings
\titleformat{\section}
  {\sffamily\bfseries\large\color{Deep}}{}{0pt}
  {}[{\vspace{1.2mm}\color{Amber}\titlerule[1.1pt]}]
\titlespacing*{\section}{0pt}{5.5mm}{2.8mm}

\titleformat{\subsection}
  {\sffamily\bfseries\normalsize\color{Teal}}{}{0pt}{}
\titlespacing*{\subsection}{0pt}{3.6mm}{1.4mm}

\setlist[itemize]{leftmargin=4.2mm, itemsep=1.1mm, topsep=1.3mm, parsep=0pt,
                  label=\textcolor{Teal}{\footnotesize$\blacktriangleright$}}

% ------------------------------------------------------------------ header/footer
\pagestyle{fancy}
\fancyhf{}
\renewcommand{\headrulewidth}{0pt}
\renewcommand{\footrulewidth}{0pt}
\fancyhead[L]{\sffamily\scriptsize\color{Slate}PM RESEARCH WATCH\;\textcolor{Rule}{\textbar}\;\DIGESTDATE}
\fancyhead[R]{\sffamily\scriptsize\color{Slate}Particulate matter monitoring \& health effects}
\fancyfoot[C]{\sffamily\scriptsize\color{Slate}\thepage}

% ------------------------------------------------------------------ components
% metric tile
\newcommand{\metric}[3]{%
  \begin{tcolorbox}[enhanced, colback=white, colframe=Rule, boxrule=0.5pt,
    arc=1.6mm, left=2.2mm, right=2.2mm, top=1.8mm, bottom=1.8mm,
    borderline west={1.7mm}{0pt}{#3}, nobeforeafter, valign=center,
    height=17mm, width=\linewidth]
    \RaggedRight
    {\sffamily\bfseries\fontsize{17}{17}\selectfont\color{#3}#1}\\[0.5mm]
    {\sffamily\fontsize{6.5}{7.4}\selectfont\color{Slate}#2}
  \end{tcolorbox}}

% pull-quote / key insight box
\newtcolorbox{keybox}[1][]{enhanced, colback=Mist, colframe=Mist,
  arc=1.6mm, left=3.2mm, right=3.2mm, top=2.4mm, bottom=2.4mm,
  borderline west={1.7mm}{0pt}{Amber}, #1}

% subtopic band
\newcommand{\band}[2]{%
  \vspace{3.2mm}\noindent
  \begin{tcolorbox}[enhanced, colback=#2, colframe=#2, arc=1.2mm,
    left=2.6mm, right=2.6mm, top=1.1mm, bottom=1.1mm, boxrule=0pt,
    width=\linewidth, nobeforeafter]
    {\sffamily\bfseries\small\color{white}#1}
  \end{tcolorbox}\vspace{1.6mm}}

% one paper entry: {short cite}{title}{journal}{doi}{body}
\newcommand{\paperentry}[5]{%
  \par\vspace{1.9mm}\noindent
  {\sffamily\bfseries\footnotesize\color{Ink}#1}\;%
  {\sffamily\fontsize{7}{8}\selectfont\color{Slate}#3}\\[0.4mm]
  {\footnotesize\itshape\color{Deep}#2}\\[0.6mm]
  {\footnotesize #4}\\[0.3mm]
  {\fontsize{6.6}{7.6}\selectfont\color{Slate}\texttt{doi:}\,\href{https://doi.org/#5}{\texttt{\detokenize{#5}}}}
  \par}

% caption style
\newcommand{\figcap}[1]{\par\vspace{1.0mm}
  {\fontsize{7.2}{8.6}\selectfont\color{Slate}\textsf{\textbf{Fig.}}\; #1\par}\vspace{1.5mm}}

\providecommand{\permil}{\textperthousand}

\setlength{\parindent}{0pt}
\setlength{\parskip}{0pt}
\renewcommand{\arraystretch}{1.18}


\begin{document}

% ==================================================================== MASTHEAD
\thispagestyle{empty}
\begin{tikzpicture}[remember picture, overlay]
  \fill[Deep] (current page.north west) rectangle
      ($(current page.north east)+(0,-67mm)$);
  \fill[Amber] ($(current page.north west)+(0,-67mm)$) rectangle
      ($(current page.north east)+(0,-68.4mm)$);
  \foreach \x/\y/\r in {22/12/0.9, 41/26/1.5, 63/9/0.7, 88/31/1.1, 112/17/1.9,
                        138/29/0.8, 159/13/1.3, 178/54/1.0, 196/21/1.6,
                        30/44/0.6, 74/57/1.2, 125/48/0.9, 168/6/0.8, 191/59/0.7,
                        52/18/1.0, 100/60/0.6, 148/20/1.1, 66/49/1.4, 205/45/0.8}
    \fill[Sky, opacity=0.30] ($(current page.north west)+(\x mm,-\y mm)$) circle (\r mm);
\end{tikzpicture}

\vspace*{4mm}
{\sffamily\fontsize{7.6}{9}\selectfont\color{Sky}%
 AUTOMATED DAILY LITERATURE DIGEST\;\textbullet\;PhD RESEARCH WATCH\par}
\vspace{2.2mm}
{\sffamily\bfseries\fontsize{27}{29}\selectfont\color{white}%
 Particulate Matter\par}
{\sffamily\fontsize{27}{29}\selectfont\color{Sky}%
 Monitoring \& Health Effects\par}
\vspace{5.0mm}
{\sffamily\fontsize{8.4}{10}\selectfont\color{white}%
 Issue 2026-09-15\;\;\textcolor{Amber}{\textbar}\;\;PubMed entry window 15 September 2026\;\;%
 \textcolor{Amber}{\textbar}\;\;22 records screened in scope\par}
\vspace{18mm}

% -------------------------------------------------------------- metric strip
\noindent
\begin{minipage}[t]{0.190\linewidth}\metric{22}{RECORDS IN SCOPE\\(19 EXCLUDED, 4 HELD)}{Deep}\end{minipage}\hfill
\begin{minipage}[t]{0.190\linewidth}\metric{9}{OF 10 SUBTOPIC\\BANDS POPULATED}{Teal}\end{minipage}\hfill
\begin{minipage}[t]{0.190\linewidth}\metric{9}{EFFECT ESTIMATES\\(ALL WITH A CI)}{Sky}\end{minipage}\hfill
\begin{minipage}[t]{0.190\linewidth}\metric{8}{RECORDS CARRYING\\NO PMID}{Amber}\end{minipage}\hfill
\begin{minipage}[t]{0.190\linewidth}\metric{5}{TIER-A\\RECORDS}{Clay}\end{minipage}

\vspace{6mm}

% -------------------------------------------------------------- signal of day
\section*{\sffamily\bfseries\color{Deep}Signal of the day}
\vspace{-1.5mm}{\color{Amber}\rule{\linewidth}{1.1pt}}\vspace{2.5mm}

\begin{keybox}
{\sffamily\bfseries\small\color{Deep}Per unit mass, wildfire PM$_{2.5}$ carries about
4.6 times the mortality hazard of ordinary urban PM$_{2.5}$ --- but the exposure contrast
behind that number is a thirtyfold difference in concentration, and that is the part to
read carefully.}\\[1.4mm]
\footnotesize
Su et al.\ (\textit{Lancet Oncol}) followed \textbf{414{,}016} SEER-Medicare patients aged
65 and over with histologically confirmed lung cancer through \textbf{1{,}211{,}490}
person-years (2008--2019), with zip-code wildfire and non-wildfire PM$_{2.5}$ from
validated models and a time-varying Cox model on 3-year moving averages. Per
1\,$\mu$g/m$^3$, the hazard of death was \textbf{HR 1.0782 (95\% CI 1.0641--1.0924)} for
wildfire PM$_{2.5}$ against \textbf{1.0171 (1.0149--1.0192)} for non-wildfire ---
non-overlapping intervals, and the sharpest population-scale statement yet that smoke mass
and urban mass are not the same exposure.\\[1.0mm]
The caution is arithmetic. Median wildfire PM$_{2.5}$ was \textbf{0.28\,$\mu$g/m$^3$}
(IQR 0.21--0.39) against \textbf{8.11} for non-wildfire, so a ``per 1\,$\mu$g/m$^3$''
wildfire coefficient is extrapolated several IQR-widths beyond anything observed, and on
the per-SD scale the ordering reverses (1.0229 versus 1.0384) because an SD of wildfire
PM$_{2.5}$ is a far smaller absolute dose. Both numbers are correct; only one is a
statement about the air people breathed. The monitoring implication survives either
reading: a network reporting mass alone cannot distinguish these two exposures.
\end{keybox}

\clearpage

% -------------------------------------------------------------- exec brief
\section*{\sffamily\bfseries\color{Deep}What changed today}
\vspace{-1.5mm}{\color{Amber}\rule{\linewidth}{1.1pt}}\vspace{2.5mm}

\footnotesize
\begin{itemize}
\item \textbf{A dense day, and for once the epidemiology carries intervals.} Twenty-two
records passed screening against sixteen on 14 September, and \textbf{nine} carry a
published exposure--response estimate with a 95\% interval. Eight of the nine are PM
estimates proper --- the first issue in some weeks whose forest plot is not carried by
heat or co-pollutant surrogates. Nine of ten subtopic bands are populated;
\textit{Burden, policy \& mitigation} is the sole empty one.

\item \textbf{The DOI gate removed three duplicates that title screening called new, and
one is a repeat offender.} The Consensus sweep returned ten low-cost-sensor calibration
hits; four matched corpus titles directly, and the DOI check then removed three more ---
a multi-city PurpleAir/OPC-N3 calibration, the Sofia calibration-transfer study, and the
Bogot\'a T640X node calibration. That last record was \emph{already removed on
14 September} for the same reason: it shipped on 3 August under a completely rewritten
title. Consensus keeps re-surfacing it and no title-normalised index will ever catch it.

\item \textbf{Two findings rest on weak arithmetic, and both are named in place.} The
Himawari-8 study of Javanese coal-plant plumes reports random forest and XGBoost scoring
\textbf{exactly 1.000} on every metric --- but the target is a deterministic threshold on
the two input variables, so the models relearn the labelling rule and the 2026--2060
projection built on them is uninterpretable. The Korean agricultural analysis lifts
precursor--PM$_{2.5}$ correlations to $r \ge 0.816$ by dividing every series by one
wind-speed-derived exchange rate, which is what dividing several series by a common
quantity does. A third belongs with them: the \textit{JECH} suicide meta-analysis pools
per-increment and high-versus-low categorical PM$_{2.5}$ contrasts under one label,
producing a 38-fold short-versus-long gap that is not a duration effect.

\item \textbf{Four in-scope records were verified and deliberately held.} The MoAE
low-cost sensor calibration framework (\textit{Atmos Environ}), a CTM--InMAP burden study
over Punjab, a SW-Europe dust and sea-salt compliance analysis (\textit{Atmos Res}) and an
office breathing-zone sensor-placement chamber study (\textit{Build Environ}) all have
resolving DOIs and no abstract in Crossref, Europe PMC, Semantic Scholar, Consensus or
Scholar Gateway --- all four same-day Elsevier deposits, the failure mode logged on
14 September. They are kept \emph{out} of \texttt{seen.json}; the MoAE paper is the
tier-A pickup. Separately, \textbf{arXiv returned HTTP 406} and an empty body, a new
failure distinct from the \texttt{submittedDate} behaviour previously logged, and
OpenAlex returned HTTP 429 again.

\item \textbf{Evidence geography is tied, not led.} China and Global / multi-region sit at
\textbf{seven records each}. Sub-Saharan Africa's two matter disproportionately --- a
Nigerian cohort and South African pneumococcal surveillance, both built on exposure
surfaces with no in-country reference network.
\end{itemize}

\clearpage

% ==================================================================== ANALYTICS
\section*{\sffamily\bfseries\color{Deep}Corpus analytics}
\vspace{-1.5mm}{\color{Amber}\rule{\linewidth}{1.1pt}}\vspace{3mm}

\noindent\includegraphics[width=\linewidth]{\FIGDIR/f1_subtopics.png}
\figcap{Subtopic assignment of the 22 in-scope records, single-label by dominant
contribution. \textit{Exposure assessment \& modelling} takes 8, twice the next cluster;
\textit{Sensing, forecasting \& instrumentation} 4; \textit{Respiratory \& allergic} 3;
\textit{Other clinical endpoints} 2; and \textit{Neuro / mental health},
\textit{Cardiovascular \& metabolic}, \textit{Reproductive \& developmental},
\textit{Mechanistic toxicology} and \textit{Occupational \& indoor} one each. Nine of the
ten standing bands are populated --- \textbf{\textit{Burden, policy \& mitigation} is the
sole empty one}.}

\vspace{2mm}
\noindent\includegraphics[width=\linewidth]{\FIGDIR/f2_design_endpoint.png}
\figcap{Left --- study architecture. Measurement campaign is the largest class at 8, ahead
of modelling and inventory at 4; observational designs split 2 cohort, 2 acute and
1 ecological, with 2 reviews and one each of experimental toxicology, sensor co-location
and chamber work. Right --- health endpoint. \textbf{Thirteen of the 22 records carry no
health endpoint}, which is the aerosol-chemistry, instrument and modelling half of the
corpus; among those that do, Respiratory leads with 3 and the remaining \textbf{six}
endpoint classes --- toxicological, reproductive, cardiovascular, neuro/cognitive, cancer
and mortality --- hold one record each.}

\vspace{2mm}
\noindent\includegraphics[width=\linewidth]{\FIGDIR/f3_metric_geography.png}
\figcap{Left --- particle metric. PM$_{2.5}$ alone is the exposure in 11 of 22 records and
PM$_{2.5}$ with PM$_{10}$ jointly in a further 5; ultrafine, optical-property and
PM$_{10}$-only work contribute one record each, so size-resolved and optical measurement
remain thin against mass. Right --- geographic provenance. \textbf{China and Global /
multi-region are tied at seven records each}; no single country dominates this issue.}

\vspace{2mm}
\noindent\includegraphics[width=\linewidth]{\FIGDIR/f6_lifecourse.png}
\figcap{Life-course windows addressed. Four of the five stages are occupied and
\textbf{only puberty and adolescence is empty}, an improvement on the three-empty issues
of the past fortnight. Working-age adulthood is nominally filled but by a wind-tunnel
instrument evaluation rather than any worker cohort, so the occupational gap is real even
where the marker is not zero.}

\vspace{2mm}
\noindent\includegraphics[width=\linewidth]{\FIGDIR/f4_heatmap.png}
\figcap{Subtopic $\times$ study-architecture cross-tabulation. \textbf{Fifteen of the 81
cells are occupied.} Density sits where measurement campaigns meet exposure assessment;
each clinical row is a single cell, which is what a one-record-per-endpoint day looks
like.}

\vspace{2mm}
\noindent\includegraphics[width=\linewidth]{\FIGDIR/f5_forest.png}
\figcap{All nine interval-bearing estimates in today's corpus on a common logarithmic
axis. \textbf{Eight are PM exposure--response estimates} --- two wildfire versus
non-wildfire mortality hazards, two all-cancer incidence hazards, one haemorrhagic-stroke
risk ratio and three suicide-outcome odds ratios for PM$_{2.5}$ and PM$_{10}$ --- with the
NO$_2$ estimate shown as a co-pollutant comparator. Interval width spans a factor of
about 300, from the SEER-Medicare non-wildfire estimate (1.0149--1.0192) to the Nigerian
top-quartile hazard (1.29--2.57); that is a difference in study size and exposure
contrast, not in precision of measurement. The two short-term suicide estimates at the
foot of the plot and the long-term estimate in second place come from the same
meta-analysis on incompatible exposure scales and should not be read against each other.}

\clearpage

% ==================================================================== DIGEST
\section*{\sffamily\bfseries\color{Deep}Paper-level digest}
\vspace{-1.5mm}{\color{Amber}\rule{\linewidth}{1.1pt}}\vspace{1mm}

{\fontsize{7.4}{9}\selectfont\color{Slate}Ordered by cluster. Each entry gives the
methodological core and the finding that matters, not an abstract paraphrase.\par}

\band{Exposure assessment \& modelling\hfill 8 records}{Amber}

\paperentry{Han et al.\ 2026}
{Source apportionment of PM$_{2.5}$ oxidative potential in metropolitan Melbourne}
{Environ Pollut}
{Published PMF source contributions linked by multiple linear regression to oxidative
potential measured on \textbf{three complementary acellular assays} (DTT, AA, DCFH);
models explain \textbf{73--86\%} of OP variance. Combustion dominates OP$_v$-DTT at
52--53\% with biomass burning above traffic; OP$_v$-AA is driven by metal-rich sources
including diesel at 20--22\%; secondary nitrate contributes 28--33\% and secondary
sulphate 10\% at the second site. The assay-dependent ranking is the result --- a
single-assay OP metric would nominate a different mitigation target. Reusing an earlier
PMF solution propagates its rotational ambiguity, and no uncertainty is quoted on the
source-resolved shares. First Southern Hemisphere estimate of its kind.}
{10.1016/j.envpol.2026.129159}

\paperentry{Liu et al.\ 2026}
{Decadal changes in PM$_{2.5}$ sources and secondary aerosol drivers, Shenzhen 2014--2024}
{J Environ Manage}
{Eleven years of speciated filter sampling with PMF apportionment and machine-learning
attribution of secondary formation drivers. Annual mean fell \textbf{37.7 to
17.3\,$\mu$g\,m$^{-3}$}. Secondary sulphate was historically dominant and declined
fastest; vehicles became the largest single contributor after 2019; \textbf{secondary
nitrate barely moved and is now the binding constraint}. NO$_x$ and VOC cuts drove nitrate
and SOA declines but rising atmospheric oxidation capacity offset them, while sulphate
tracked oxidation and regional transport over local SO$_2$. Attributing a decadal trend to
oxidation capacity from a correlational model at one site cannot exclude meteorological
confounding, and the PMF factor series carry no uncertainty bounds.}
{10.1016/j.jenvman.2026.130913}

\paperentry{Zhang et al.\ 2026}
{Surface-based temperature inversion climatology across China and its relation to PM$_{2.5}$}
{Environ Pollut}
{Hourly microwave-radiometer profiles from \textbf{33 stations} over 2024--2025, validated
against radiosondes at 25 sites, used to build a surface-based inversion climatology and
relate it to PM$_{2.5}$ at hotspot stations. Winter SBI occurrence \textbf{39.2\%} against
\textbf{12.1\%} in summer, with a northwest-to-southeast gradient peaking in the Tarim
Basin and Hexi Corridor. SHAP and Spearman analysis place boundary-layer variables in the
top predictor tier monthly and daily, with effective mixing-layer height third among all
hourly predictors in every season. Two years is short for a climatology, the network is
sparse relative to the terrain it spans, and SHAP importance over correlated
meteorological predictors does not establish the causal pathway it is read as
establishing.}
{10.1016/j.envpol.2026.129172}

\paperentry{Kim et al.\ 2026}
{LASSO-based reduced-form CMAQ for ozone and PM$_{2.5}$ response to emission change}
{PLoS One}
{Surrogate for CMAQ fitted by LASSO over \textbf{118 emission scenarios} spanning seven
sectors and 17 provincial regions, with an adaptive logit transform to respect
concentration bounds. Trained on 100 simulations and tested on 18: mean RMSE
\textbf{0.1\,ppb} for ozone and \textbf{0.1\,$\mu$g\,m$^{-3}$} for PM$_{2.5}$ against the
parent model. The sparse set of retained sector--region terms is the practical output,
making the control-design problem interpretable rather than opaque. Error is quoted
against CMAQ and not against observations, so every CMAQ bias passes through undiagnosed;
118 scenarios is thin for a 119-dimensional predictor space and no out-of-envelope
extrapolation test is reported.}
{10.1371/journal.pone.0347073}

\paperentry{Aji et al.\ 2026}
{Fine-mode combustion aerosol events near Indonesian coal plants from AE-filtered Himawari-8}
{Environ Monit Assess}
{\textbf{41{,}119} daily Himawari-8 observations around 13 Java coal plants, 2015--2025,
with events defined as \r{A}ngstr\"om exponent $>1.2$ and AOD $>0.4$ (8.52\% of days).
Suralaya has the highest mean AOD at 1.01 with 51.4\% very-high-pollution days; Pacitan
the lowest at 0.46. Random forest and XGBoost return accuracy, precision, recall, F1 and
AUC-ROC all \textbf{exactly 1.000} --- on a target that is a deterministic threshold
function of the two inputs, so the models relearn the labelling rule rather than predict
anything, and the 2026--2060 Monte Carlo projection built on them is uninterpretable.
Column AOD without surface validation cannot separate plant plume from boundary-layer
aerosol.}
{10.1007/s10661-026-15804-1}

\paperentry{Ilies et al.\ 2026}
{Urban PM$_{10}$ and population exposure from satellite, meteorological and socio-spatial data}
{Front Public Health}
{MODIS-MAIAC AOD combined with four ground stations and meteorology in a regression for
PM$_{10}$ over Oradea, Romania, then folded with population density, built form, NDVI and
elevation into a composite exposure index. Cold-season maxima reach
\textbf{41.8\,$\mu$g\,m$^{-3}$} against an annual mean near 20.5 --- above WHO, inside EU
limits --- and \textbf{20.1\%} of the metropolitan population falls in
high-predisposition areas. The AOD--PM$_{10}$ relation alone is weak
(\textbf{R = 0.28}) and meteorology carries the model, which is the authors' own admission
and the standing reason satellite-only surfaces underperform at city scale. Four stations
cannot validate an intra-urban field, and the composite index is unweighted and never
checked against measured personal exposure.}
{10.3389/fpubh.2026.1940025}

\paperentry{Kim \& Park 2026}
{Loading-rate-based analysis of PM$_{2.5}$ and precursor gases in Korean agricultural areas}
{Atmosphere}
{547 days of hourly data from eight agricultural sites, comparing concentration-based with
loading-rate-based interpretation, where loading rate is concentration scaled by a
wind-speed-derived air exchange rate. Concentration correlations among PM$_{2.5}$,
NH$_3$, NO$_2$ and SO$_2$ are weak and inconsistent; loading-rate correlations are strong
and uniform at \textbf{$r \ge 0.816$}, and the diurnal shape flips from bimodal to
unimodal. The authors concede the improvement may reflect the shared wind-speed term ---
it almost certainly does, since dividing every series by one common quantity induces
correlation by construction. Recorded as a method to treat with suspicion: any dense
open-field network normalising by a common meteorological variable manufactures the same
result.}
{10.3390/atmos17090898}

\paperentry{Ahmad et al.\ 2026}
{Carbonaceous composition and multi-wavelength optical properties of PM in Lahore}
{Atmosphere}
{Six-week urban campaign in Lahore, 30 quartz filters collected, 25 usable for gravimetry
and optics, \textbf{only the first 15 analysed for OC/EC}. PM$_{2.5}$ ranged
\textbf{40--417\,$\mu$g\,m$^{-3}$} with a mean of \textbf{137 $\pm$ 84}. Organic carbon
was 76.6\% of total carbon against 23.4\% elemental. Mean absorption \r{A}ngstr\"om
exponent \textbf{1.17 $\pm$ 0.30}, with 56\% of samples in 1.0--1.5, 32\% below 1.0 and
12\% above 1.5. The authors are appropriately careful that AAE intervals describe spectral
behaviour rather than attribute sources. Fifteen OC/EC filters in one pre-monsoon season
cannot characterise Lahore's annual cycle, and the gravimetric series has no co-located
reference for traceability. Directly relevant to optical low-cost sensors, whose response
depends on exactly this absorption behaviour.}
{10.3390/atmos17090899}

\band{Sensing, forecasting \& instrumentation\hfill 4 records}{Amber}

\paperentry{Vance et al.\ 2026}
{Evaluation of the Partector 2 Pro for ultrafine particles from 3D printing}
{J Occup Environ Hyg}
{Wind-tunnel comparison of a portable diffusion-charging monitor against SMPS and CPC
during controlled printing of PLA, PETG, TPU and nylon. Number-concentration agreement is
strong across all four filaments --- regression slopes \textbf{0.88--1.05}, Pearson
$r > 0.96$, mean deviations \textbf{$-$9.6\% to $+$6.2\%}, inside the $\pm$30\%
manufacturer specification. Sizing is systematically biased low (slopes 0.82--0.90) and
worst at the fine end: \textbf{$-$27.0\% to $+$7.5\% at D10} and $-$28.3\% to $-$7.3\%
across D25--D75. That pattern is the operational point for any distributed ultrafine
network --- the instrument is fit for concentration surveillance but its size distribution
must not be used to infer deposition fraction. Wind tunnel only, with no humidity or
co-pollutant interference testing.}
{10.1080/15459624.2026.2724020}

\paperentry{Wathore et al.\ 2026}
{Environmental and temporal performance of ML calibration for a low-cost PM$_{2.5}$ sensor}
{Aerosol Air Qual Res}
{Long-term co-location in Vishakhapatnam comparing Random Forest and XGBoost against
linear regression, adding temperature, relative humidity, hour-of-day and month-of-year
stepwise, plus a hybrid ensemble. The best single model cuts RMSE \textbf{58\%} against
the linear baseline and the hybrid reaches \textbf{63\%}, meeting US EPA low-cost
PM$_{2.5}$ performance criteria. The contribution is the evaluation design rather than the
models: performance is reported binned across environmental and temporal strata and across
AQI categories, which exposes where a calibration degrades instead of burying it in a
pooled $R^2$. No absolute RMSE in $\mu$g\,m$^{-3}$ is given, and single-site training
leaves transferability untested --- the same limitation already flagged for the Sofia
calibration-transfer work in this corpus.
\textit{Connector-sourced; outside the 15 September entry window.}}
{10.1007/s44408-026-00094-6}

\paperentry{Komiljon et al.\ 2026}
{Non-cumulative, size-specific calibration of low-cost PM sensors under simulated drilling}
{Atmosphere}
{Controlled laboratory evaluation of PMS5003 and Sniffer4D under simulated concrete-slab
drilling, with six correction models --- linear, polynomial, Random Forest, XGBoost, ANN,
Kalman --- fitted independently per size fraction. Random Forest and ANN perform best,
\textbf{$R^2 > 0.89$ for PM$_1$} and \textbf{$> 0.87$ for PM$_{2.5}$} at 50\,s averaging.
The premise is the useful part: reporting non-cumulative rather than cumulative bins,
since cumulative output conflates fractions with different deposition behaviour and
different regulatory meaning. Against that, a single drilling source in a chamber produces
a mineral-dust size distribution that will not transfer to combustion aerosol, and
PM$_{10}$ performance is not reported alongside the finer fractions.
\textit{Connector-sourced; Crossref resolves only the preprint of record.}}
{10.20944/preprints202604.1260.v1}

\paperentry{Wang et al.\ 2026}
{14-day PM$_{2.5}$ and O$_3$ forecasting over eastern China with a CNN--Transformer}
{Environ Sci Technol}
{Hybrid architecture combining multilevel regional meteorological fields with emission
inventories at a 14-day lead. Relative $R^2$ gains over baseline are \textbf{40\% for
O$_3$} (0.58 to 0.81) and \textbf{18\% for PM$_{2.5}$} (0.44 to 0.52); 850\,hPa fields add
5\% for O$_3$ and emissions add 9.8\% for PM$_{2.5}$; 66--83\% of O$_3$ episodes are
captured in the north. Read honestly, PM$_{2.5}$ skill remains poor --- $R^2$ 0.52 at
14 days is not decision-grade, and the paper's own framing concedes O$_3$ is the stronger
case. No episode-capture rate is reported for PM$_{2.5}$, which is the number a forecast
user would act on, and gridded reanalysis truth hides station-scale error.}
{10.1021/acs.est.6c05923}

\clearpage

\band{Respiratory \& allergic\hfill 3 records}{Sage}

\paperentry{Kremer et al.\ 2026}
{Pneumococcal population structure modifies the effect of air pollution on invasive disease}
{Nat Microbiol}
{Nineteen years of South African national surveillance --- \textbf{59{,}017} invasive
pneumococcal disease cases with \textbf{4{,}350} sequenced isolates --- in Bayesian
spatiotemporal models linking environment to pathogen diversity. PM$_{2.5}$ is positively
associated with IPD, but \textbf{the lag structure varies by serotype}: only serotypes 4,
8 and 23F confer immediate risk, and high prevalence of lineage GPSC21 (serotype 19F)
shifts the population-level response to immediate. Cumulative 8-week effects for moderate
humidity and cold minimum temperature are $+5\%$ and $+4\%$. No confidence interval
accompanies the headline pollution association and the serotype-stratified lag models
present a large multiple-comparison surface. Still the first demonstration that pathogen
genotype modifies an air-quality exposure--response, and a direct argument for pairing
dense PM monitoring with pathogen surveillance.}
{10.1038/s41564-026-02458-5}

\paperentry{Zhu et al.\ 2026}
{Extreme cold events and paediatric asthma visits, modified by ambient CO and SO$_2$}
{J Asthma Allergy}
{\textbf{62{,}167} paediatric asthma outpatient visits at one tertiary hospital in Nanjing,
2019--2024, quasi-Poisson DLNM over lag 0--14. Extreme cold events show \textbf{no overall
association}, but cumulative RR reaches \textbf{2.60 (95\% CI 1.13--5.95)} in the high-CO
stratum and \textbf{3.41 (1.15--10.11)} in the high-SO$_2$ stratum. The intervals are very
wide and the strata are subsets of an already null main effect, so this is
hypothesis-generating; notably PM$_{2.5}$ and PM$_{10}$ did not emerge as modifiers.
Single-hospital catchment with city-average exposure. Relevant here mainly as a reminder
that pollutant-by-meteorology interaction terms demand far denser exposure data than one
municipal station provides.}
{10.2147/jaa.s483101}

\paperentry{Martinez \& Nadeau 2026}
{Early-life exposure to air pollutants and childhood asthma}
{Curr Allergy Asthma Rep}
{Narrative synthesis of recent cohort and mechanistic work on the prenatal and early
postnatal window. Reports epigenetic modification in airway epithelium, a convergence
point shared between tobacco smoke exposure and genetic asthma predisposition, new serum
biomarkers of exposure, and a murine result in which asthma-predisposing information
passes from exposed dam to offspring \textbf{via NK cells}. No pooled estimate, no effect
size, and no stated search strategy or inclusion criteria, so study selection is
unauditable and the piece cannot be weighted against the systematic reviews in this
window. Useful as a pointer to the NK-cell transmission mechanism and the biomarker
candidates, and for nothing quantitative.}
{10.1007/s11882-026-01293-2}

\band{Other clinical endpoints\hfill 2 records}{Slate}

\paperentry{Su et al.\ 2026}
{Wildfire-related PM$_{2.5}$ and overall survival of older lung cancer patients}
{Lancet Oncol}
{\textbf{414{,}016} SEER-Medicare patients aged 65+, \textbf{1{,}211{,}490} person-years,
2008--2019, time-varying Cox on 3-year moving averages of zip-code wildfire and
non-wildfire PM$_{2.5}$. Per 1\,$\mu$g\,m$^{-3}$: \textbf{HR 1.0782 (95\% CI
1.0641--1.0924)} for wildfire against \textbf{1.0171 (1.0149--1.0192)} for non-wildfire,
a 4.6-fold difference in per-unit potency. The exposure contrast is the weakness --- median
wildfire PM$_{2.5}$ was 0.28\,$\mu$g\,m$^{-3}$ (IQR 0.21--0.39) against 8.11 for
non-wildfire, so the per-unit hazard is extrapolated far outside the observed range, and
the per-SD estimates (1.0229 versus 1.0384) reverse the ordering. Zip-code assignment with
no residential history. Species-resolved rather than mass-only monitoring is the direct
implication.}
{10.1016/s1470-2045(26)00295-0}

\paperentry{Mtenga et al.\ 2026}
{Ambient PM$_{2.5}$ exposure and cancer incidence in Nigerian adults}
{Cancer Epidemiol}
{First prospective sub-Saharan African cohort on this question. 5{,}621 adults, median
follow-up 4.7 years, 147 incident cancers confirmed histopathologically, mean PM$_{2.5}$
\textbf{50.4\,$\mu$g\,m$^{-3}$} --- tenfold the WHO guideline. Each 10\,$\mu$g\,m$^{-3}$
gives \textbf{HR 1.12 (95\% CI 1.06--1.19)} for all-cancer incidence and the top quartile
\textbf{1.82 (1.29--2.57)}; mediation analysis attributes \textbf{41.3\%} to oxidative
stress and inflammation. Site-specific signals for lung, breast and haematopoietic cancer
rest on very small strata at 147 total events, and biomarkers measured once in a subsample
cannot establish the temporal ordering mediation requires. The 1\,km satellite--land-use
regression exposure surface has no ground reference network to anchor it.}
{10.1016/j.canep.2026.103232}

\clearpage

\band{Neuro / mental health\hfill 1 record}{Deep}

\paperentry{Pearson et al.\ 2026}
{Environmental exposure and risk of suicide, attempted suicide and ideation}
{J Epidemiol Community Health}
{Forty-five studies screened across air, noise, light and water pollution, 29 pooled.
Short-term PM$_{2.5}$ \textbf{OR 1.0139 (95\% CI 1.0001--1.0279)} over 13 studies and
PM$_{10}$ \textbf{1.0100 (1.0007--1.0195)} over 7; long-term PM$_{2.5}$ \textbf{1.5254
(1.1351--2.0481)} over 6 and NO$_2$ \textbf{1.3173 (1.0641--1.6308)}. The short-term
PM$_{2.5}$ lower bound sits at 1.0001, so significance is nominal. More seriously, the
\textbf{38-fold gap between short- and long-term PM$_{2.5}$ point estimates is an artefact
of pooling per-increment against high-versus-low categorical contrasts under one label};
the two are not on a common scale and must not be read as a duration--response.
Newcastle-Ottawa assessment only, with no weighting for exposure misclassification.}
{10.1136/jech-2025-225491}

\band{Cardiovascular \& metabolic\hfill 1 record}{Teal}

\paperentry{Xu et al.\ 2026}
{Ambient fine particulate matter and haemorrhagic stroke in Yancheng, China}
{Front Public Health}
{Individual-level time-stratified case-crossover on \textbf{2{,}312} haemorrhagic stroke
admissions across four Yancheng hospitals, 2022--2024, conditional logistic regression.
Strongest at lag01: each 10\,$\mu$g\,m$^{-3}$ of PM$_{2.5}$ associated with a
\textbf{3.24\% increase (95\% CI 0.55--6.00)}, with lag0 2.50\% (0.11--4.95) and lag1
2.76\% (0.32--5.27). The association survives adjustment for O$_3$ but \textbf{loses
significance once SO$_2$ or NO$_2$ enters} --- which is the finding that matters, because
it means the PM$_{2.5}$ coefficient is not separable from co-emitted combustion products
at this station density. City-average exposure assigned to residence leaves within-city
gradients unmeasured, and a four-hospital catchment is not the full case population.}
{10.3389/fpubh.2026.1882030}

\band{Reproductive \& developmental\hfill 1 record}{Sky}

\paperentry{Sun et al.\ 2026}
{Air pollutant mixtures and endometrial receptivity in the peri-implantation window}
{Environ Res}
{\textbf{1{,}349} women across 14 cities in Liaoning, with quantile g-computation and BKMR
on a six-pollutant mixture against sonographic endometrial thickness and volume. A
simultaneous one-quartile increase in the mixture is associated with \textbf{0.5\,mm} less
thickness and \textbf{0.28\,ml} less volume, with CO and O$_3$ carrying the largest
weights and PM$_{2.5}$ and PM$_{10}$ not dominating. \textbf{No confidence intervals are
reported for the mixture effects}, which for a BKMR result is a substantive omission. A
0.5\,mm shift is at or below the reproducibility of transvaginal endometrial measurement,
so the clinical meaning is unclear even if the association holds. Monthly gridded exposure
with no personal or indoor component, against an implantation window measured in days.}
{10.1016/j.envres.2026.125682}

\band{Mechanistic toxicology\hfill 1 record}{Violet}

\paperentry{Garner et al.\ 2026}
{Gestational nanoparticle inhalation disrupts placental zone structure in rats}
{Reprod Toxicol}
{Pregnant Sprague Dawley rats exposed by whole-body inhalation to nano-TiO$_2$ at
\textbf{9.23 $\pm$ 0.39\,mg\,m$^{-3}$} from gestational day 5 to 19, placentas harvested
GD20 for histology. Placental weight and placental efficiency were \textbf{unchanged}, but
decidua and labyrinth zone sizes fell, with compensatory increases in blood space number
and maternal blood space area. The compensation is the interesting part: structural
remodelling without a mass-level phenotype means placental weight is a poor screening
endpoint for ultrafine exposure. The exposure is roughly three orders of magnitude above
ambient ultrafine mass and engineered TiO$_2$ is a poor surrogate for combustion
ultrafines in surface chemistry, so this bounds mechanism and not risk. No litter-level
statistics are reported.}
{10.1016/j.reprotox.2026.109361}

\band{Occupational \& indoor\hfill 1 record}{Clay}

\paperentry{Chen et al.\ 2026}
{Indoor air protection under extreme air pollution risk: evidence from a sandstorm in China}
{Indoor Air}
{Difference-in-differences on Baidu search intensity for air purifiers and indoor
air-quality information around China's most intense sandstorm in a decade (early 2021).
Affected regions show roughly \textbf{18\%} higher search intensity than unaffected ones,
rising to about \textbf{33\%} in historically cleaner areas and about \textbf{200\%} where
baseline purifier sales were lowest. The gradient is the useful result: protective
response scales inversely with prior exposure experience and prior equipment stock. Search
volume proxies intent, not purchase, use, or any realised reduction in indoor PM$_{2.5}$
--- no concentration data enter the study at all --- and regional treatment assignment with
no individual exposure makes the effect ecological throughout.}
{10.1155/ina/2274208}

\clearpage

% ==================================================================== TABLE
\section*{\sffamily\bfseries\color{Deep}Corpus register}
\vspace{-1.5mm}{\color{Amber}\rule{\linewidth}{1.1pt}}\vspace{2.5mm}

{\fontsize{6.7}{8.4}\selectfont
\setlength{\tabcolsep}{2.4pt}
\begin{longtable}{@{}>{\RaggedRight}p{25mm}>{\RaggedRight}p{41mm}>{\RaggedRight}p{27mm}%
>{\RaggedRight}p{23mm}>{\RaggedRight}p{18mm}>{\RaggedRight}p{15mm}@{}}
\toprule
\textsf{\textbf{First author}} & \textsf{\textbf{Focus}} & \textsf{\textbf{Journal}} &
\textsf{\textbf{Design}} & \textsf{\textbf{Metric}} & \textsf{\textbf{Region}}\\
\midrule
\endfirsthead
\toprule
\textsf{\textbf{First author}} & \textsf{\textbf{Focus}} & \textsf{\textbf{Journal}} &
\textsf{\textbf{Design}} & \textsf{\textbf{Metric}} & \textsf{\textbf{Region}}\\
\midrule
\endhead
\multicolumn{6}{@{}l@{}}{\cellcolor{Mist}\textsf{\textbf{\textcolor{Deep}{Neuro / mental health}}}}\\[0.6mm]
\href{https://doi.org/10.1001/jamanetworkopen.2026.25015}{Dales} & Wildfire-sourced PM$_{2.5}$ and ED visits for migraine and other primary headache & \textit{JAMA Netw Open} & Case-crossover & PM$_{2.5}$ (wildfire) & Canada\\
\href{https://doi.org/10.1093/ageing/afag215}{Glover} & Umbrella meta-review of urban environmental exposures and cognitive health & \textit{Age Ageing} & Umbrella review / meta-analysis & PM$_{2.5}$ & Global\\
\href{https://doi.org/10.1016/j.nsa.2026.107019}{Goossens} & In utero air pollution exposure and pre-/postnatal brain development & \textit{Neurosci Appl} & Narrative review & PM$_{2.5}$ / PM$_{10}$ / PM$_{\mathrm{coarse}}$ & Global\\
\href{https://doi.org/10.1016/j.isci.2026.116617}{Ma} & PM$_{2.5}$-induced neurotoxicity and depression: exposure risks and mechanisms & \textit{iScience} & Narrative review & PM$_{2.5}$ + components & Global\\
\href{https://doi.org/10.3760/cma.j.cn112150-20250804-00752}{Pan} & Maternal PM exposure and offspring neurodevelopment: effects and mechanisms & \textit{Zhonghua Yu Fang Yi Xue Za Zhi} & Narrative review & PM (bulk) & China\\
\addlinespace[0.9mm]
\multicolumn{6}{@{}l@{}}{\cellcolor{Mist}\textsf{\textbf{\textcolor{Teal}{Cardiovascular \& metabolic}}}}\\[0.6mm]
\href{https://doi.org/10.3760/cma.j.cn112338-20251120-00843}{Huang} & PM$_{2.5}$ and dyslipidaemia in rural Chinese children; dietary effect modification & \textit{Zhonghua Liu Xing Bing Xue Za Zhi} & Cross-sectional (multistage) & PM$_{2.5}$ & China\\
\href{https://doi.org/10.32394/pe/218966}{Kazmierska} & Air pollution and cardiovascular disease: review of current data & \textit{Przegl Epidemiol} & Narrative review & PM$_{2.5}$ / PM$_{10}$ & Global\\
\href{https://doi.org/10.1007/s00125-026-06814-2}{Xi} & Air pollution, built environment and temperature vs. prediabetes phenotypes & \textit{Diabetologia} & Prospective cohort & PM$_{2.5}$ & Germany\\
\href{https://doi.org/10.1161/JAHA.125.046917}{Yang} & Short-term PM$_{2.5}$ and hospitalisation/ED use in atrial fibrillation & \textit{J Am Heart Assoc} & Case-crossover & PM$_{2.5}$ (wildfire-influenced) & USA\\
\addlinespace[0.9mm]
\multicolumn{6}{@{}l@{}}{\cellcolor{Mist}\textsf{\textbf{\textcolor{Sky}{Reproductive \& developmental}}}}\\[0.6mm]
\href{https://doi.org/10.1016/j.jes.2025.10.043}{Chen} & Prenatal air pollution, fetal growth trajectory phenotypes and child neurodevelopment & \textit{J Environ Sci (China)} & Prospective cohort & PM$_{2.5}$ / PM$_{10}$ & China\\
\href{https://doi.org/10.1016/j.ecoenv.2026.120566}{Huang} & Joint prenatal bisphenol + air-pollutant mixture exposure and preschool IQ & \textit{Ecotoxicol Environ Saf} & Prospective cohort (mixtures) & PM$_{2.5}$ / PM$_{10}$ & China\\
\href{https://doi.org/10.1007/s40572-026-00552-8}{Peterson} & Ambient air pollution, PAHs and pubertal development: critical review & \textit{Curr Environ Health Rep} & Critical review & PM$_{2.5}$ components & Global\\
\href{https://doi.org/10.1016/j.envres.2026.125313}{Xie} & Long-term PM$_{2.5}$ and timing of pubertal development in boys and girls & \textit{Environ Res} & Prospective cohort & PM$_{2.5}$ & China\\
\addlinespace[0.9mm]
\multicolumn{6}{@{}l@{}}{\cellcolor{Mist}\textsf{\textbf{\textcolor{Sage}{Respiratory \& allergic}}}}\\[0.6mm]
\href{https://doi.org/10.1007/s00484-026-03289-z}{Bayar Muluk} & Climate change and otolaryngology: airway disease impacts and emission reduction & \textit{Int J Biometeorol} & Narrative review & PM$_{2.5}$ & Global\\
\href{https://doi.org/10.4103/mgr.MEDGASRES-D-25-00164}{Liao} & Short-term air pollution, respiratory symptoms and lung function in asthma & \textit{Med Gas Res} & Panel study & PM$_{2.5}$ / PM$_{10}$ & China\\
\href{https://doi.org/10.3390/epidemiologia7040097}{Sanchez de Prada} & Ambient PM and asthma-coded admissions in children 0-3 y, Spain & \textit{Epidemiologia} & Ecological panel & PM$_{2.5}$ / PM$_{10}$ & Spain\\
\href{https://doi.org/10.1016/j.jhazmat.2026.143076}{Yang} & Allergenic vs. toxic risk of airborne seed hairs compared with ambient PM$_{2.5}$ & \textit{J Hazard Mater} & Field measurement & PM$_{2.5}$ + bioaerosol & China\\
\addlinespace[0.9mm]
\multicolumn{6}{@{}l@{}}{\cellcolor{Mist}\textsf{\textbf{\textcolor{Violet}{Mechanistic toxicology}}}}\\[0.6mm]
\href{https://doi.org/10.1016/j.jes.2025.12.074}{Ding} & PM$_{2.5}$ promotes vascular calcification via miR-27b-5p/HES1 signalling & \textit{J Environ Sci (China)} & In vitro + organoid + murine & PM$_{2.5}$ & China\\
\href{https://doi.org/10.1007/s00360-026-01698-5}{Maraschi} & Osmoregulatory disruption and metal bioaccumulation from settleable atmospheric PM & \textit{J Comp Physiol B} & Controlled exposure (ecotox) & Settleable APM & Brazil\\
\href{https://doi.org/10.1016/j.jes.2025.11.018}{Wu} & Macrophage NLRP3 inflammasome in real-ambient PM-induced hepatic glucose dysregulation & \textit{J Environ Sci (China)} & Real-ambient murine + co-culture & Real-ambient PM & China\\
\addlinespace[0.9mm]
\multicolumn{6}{@{}l@{}}{\cellcolor{Mist}\textsf{\textbf{\textcolor{Clay}{Occupational \& indoor}}}}\\[0.6mm]
\href{https://doi.org/10.3389/fmedt.2026.1796828}{Chen} & Nested suction instrument to cut intraoperative surgical-smoke PM$_{2.5}$ & \textit{Front Med Technol} & Bench / device experiment & PM$_{2.5}$ (surgical smoke) & China\\
\href{https://doi.org/10.1093/annweh/wxag055}{Cole} & Determinants of respirable crystalline silica exposure in Australian tunnelling & \textit{Ann Work Expo Health} & Observational exposure study & Respirable dust / RCS & Australia\\
\href{https://doi.org/10.1016/j.jhazmat.2026.143081}{Ding} & Stage-resolved hair Zn-isotope dynamics as a firefighting combustion-exposure tracer & \textit{J Hazard Mater} & Repeated-measures biomarker & Combustion-derived PM & China\\
\href{https://doi.org/10.3390/jox16040119}{Ramsden} & Xenobiotic hazards in aircraft cabin air & \textit{J Xenobiot} & Review + risk economics & Ultrafine particles & Global\\
\addlinespace[0.9mm]
\multicolumn{6}{@{}l@{}}{\cellcolor{Mist}\textsf{\textbf{\textcolor{Amber}{Sensing, forecasting \& instrumentation}}}}\\[0.6mm]
\href{https://doi.org/10.1016/j.jenvman.2026.130515}{Shi} & Adaptive spatiotemporal graph transformer for multi-site multi-step PM$_{2.5}$ forecasting & \textit{J Environ Manage} & Deep-learning model & PM$_{2.5}$ & China / India\\
\href{https://doi.org/10.3390/bios16070380}{Tsuchiya} & Seasonal and indoor-outdoor characteristics of fluorescent aerosol particles in offices & \textit{Biosensors} & Real-time sensor campaign & Bioaerosol / coarse & Japan\\
\addlinespace[0.9mm]
\multicolumn{6}{@{}l@{}}{\cellcolor{Mist}\textsf{\textbf{\textcolor{Amber}{Exposure assessment \& modelling}}}}\\[0.6mm]
\href{https://doi.org/10.1029/2025GH001621}{Solorzano-Araque} & Multiscale exposure modelling: LOTOS-EUROS + mobility surveys + particle toxicity & \textit{GeoHealth} & CTM + mobility integration & PM$_{2.5}$ / PM$_{10}$ & Colombia\\
\addlinespace[0.9mm]
\multicolumn{6}{@{}l@{}}{\cellcolor{Mist}\textsf{\textbf{\textcolor{Coral}{Burden, policy \& mitigation}}}}\\[0.6mm]
\href{https://doi.org/10.1002/hsr2.72886}{Karamzad} & Ischaemic heart disease burden in MENA 1990-2021; PM attributable fraction 28.2\% & \textit{Health Sci Rep} & GBD secondary analysis & Ambient PM & MENA\\
\href{https://doi.org/10.1371/journal.pone.0347989}{Rees} & HEPA filtration to reduce respiratory infection in care homes & \textit{PLoS One} & RCT process evaluation & Filtered PM & UK\\
\href{https://doi.org/10.1016/j.jes.2026.01.009}{Shan} & Policy-driven emission reductions and post-pandemic rebound, high-resolution inventories & \textit{J Environ Sci (China)} & Emission inventory & PM emissions & China\\
\addlinespace[0.9mm]
\multicolumn{6}{@{}l@{}}{\cellcolor{Mist}\textsf{\textbf{\textcolor{Slate}{Other clinical endpoints}}}}\\[0.6mm]
\href{https://doi.org/10.3390/dj14070440}{Lupita} & PM$_{2.5}$, preventive-care capacity and caries/edentulism burden across the EU & \textit{Dent J} & Ecological panel & PM$_{2.5}$ & EU-27\\
\href{https://doi.org/10.1016/j.jdin.2026.06.008}{Park} & Long-term PM exposure and risk of androgenetic alopecia & \textit{JAAD Int} & Nationwide cohort & PM$_{2.5}$ / PM$_{10}$ & South Korea\\
\href{https://doi.org/10.3171/2026.3.PEDS2615}{Sun} & Environmental exposure and area deprivation vs. paediatric hydrocephalus outcomes & \textit{J Neurosurg Pediatr} & Registry cohort + geospatial & PM$_{2.5}$ (satellite-derived) & USA\\
\addlinespace[0.9mm]
\bottomrule
\end{longtable}}


\vspace{4mm}
\begin{keybox}
{\sffamily\bfseries\footnotesize\color{Deep}Provenance and method}\\[1.2mm]
{\fontsize{7.2}{8.8}\selectfont
Entry window \textbf{2026-09-15 $\to$ 2026-09-15}, contiguous with the 14 September issue;
no gap. Built on 16 September under the standing rule that the current day is not opened
before 22:00 local.
\textbf{Sources.} PubMed \texttt{[Date - Entry]}, run as two separate queries --- health
axis and instrumentation axis --- through both the local harvester and the connector:
\textbf{30 unique candidates}, the connector contributing 4 the harvester missed and the
harvester 7 the connector missed, with the instrumentation query returning a single record
already inside the health set. Europe PMC \texttt{CREATION\_DATE} 40, predominantly
off-topic MDPI deposits. Crossref windowed by ISSN across the tracked non-PubMed journals:
\textbf{57 deposits}. \textbf{arXiv returned HTTP 406 Not Acceptable} and an empty body ---
a new failure mode, distinct from the behaviour previously logged for
\texttt{submittedDate}. OpenAlex returned HTTP 429 again, as on every recent run. Consensus
sweep on low-cost sensor calibration returned 10.
\textbf{Screening.} 130 raw records, 120 after cross-source dedupe, 111 unseen.
\textbf{22 shipped, 23 excluded or held} --- 7 duplicates, of which \textbf{3 were caught
by the DOI gate after title screening called them new}; 12 out of scope; and \textbf{4
held} whose DOIs resolve but for which no abstract exists in Crossref, Europe PMC,
Semantic Scholar, Consensus or Scholar Gateway. The held records are deliberately kept out
of \texttt{seen.json} so a later issue can carry them.
\textbf{Trial watch.} One registry record updated in window --- NCT07713264 (VISIONS,
Stanford University, outdoor workers, interventional, $n=50$). \texttt{analyze\_endpoints}
was run per protocol: all three registered endpoints are psychosocial scales and
\textbf{no air-quality or PM measure exists in the protocol}, the second consecutive
topical hit whose PM framing is not carried by any measured PM variable. Accumulated to
\texttt{trials.json} and surfaced in the weekly rollup, not here.
\textbf{Labels.} Subtopic, design, particle-metric and geography are assigned from title,
abstract, keywords and MeSH; assignment is single-label by dominant contribution and is a
judgement, not a controlled vocabulary. One new design label
(\textit{Bayesian spatiotemporal model}) was added to the classifier in the same edit that
wrote it onto a record; every other label reused the archive.
Effect estimates are transcribed verbatim from abstracts and none is rescaled, with the
single exception of the Yancheng percent-change, converted to a risk ratio and flagged as
such.
Content derived from PubMed and Crossref metadata; abstracts remain the copyright of their
respective publishers.\par}
\end{keybox}

\end{document}
